\documentclass[10pt,a4paper]{article}

\usepackage[utf8]{inputenc}
\usepackage[T1]{fontenc}
\usepackage{lmodern}
\usepackage[margin=0.6in]{geometry}
\usepackage{enumitem}
\usepackage{titlesec}
\usepackage{amsmath,amssymb}
\usepackage{xcolor}
\usepackage[hidelinks]{hyperref}

\hypersetup{
    colorlinks=true,
    linkcolor=blue!60!black,
    urlcolor=blue!60!black
}

\pagestyle{empty}
\setlength{\parindent}{0pt}
\setlist[itemize]{leftmargin=1.5em, itemsep=2pt, topsep=3pt, parsep=0pt}

\titleformat{\section}
  {\large\bfseries\uppercase}
  {}{0em}{}
  [\titlerule]
\titlespacing*{\section}{0pt}{10pt}{6pt}

\newcommand{\entryheader}[4]{
  \textbf{#1} \hfill \textit{#2} \\
  \textit{#3} \hfill #4 \vspace{2pt}
}

\newcommand{\projectheader}[3]{
  \textbf{#1} $\vert{}$ \textit{#2} \hfill \textit{#3} \vspace{2pt}
}

\begin{document}

% -------------------- HEADER --------------------
\begin{center}
    {\LARGE \textbf{Pranav Joshi}} \\[4pt]
    \href{mailto:pranav.joshi@alumni.iitgn.ac.in}{pranav.joshi@alumni.iitgn.ac.in} $\vert{}$
    \href{mailto:pranavajoshi04@gmail.com}{pranavajoshi04@gmail.com} $\vert{}$
    +91 8698955551 \\[2pt]
    \href{https://github.com/pranav-joshi-iitgn}{GitHub} $\vert{}$
    \href{https://www.linkedin.com/in/pranav-joshi-77487b232}{LinkedIn} $\vert{}$
    \href{https://pranavjoshiiitgn.notion.site/54b8560288a346d39574e2b4cd6959ad?v=12b71196f4794d6485295473ab547e0f}{Academic Portfolio} $\vert{}$
    \href{https://medium.com/@pranav.joshi_16395}{Medium}
\end{center}

\vspace{2pt}
Engineer and researcher specializing in FPGA hardware acceleration, ultrasonic image reconstruction, real-time signal processing, and statistical modeling. 
Proven expertise in deriving and implementing mathematical and hardware-software co-design frameworks; from high-throughput beamforming architectures to physiological time-series modeling.

% -------------------- EDUCATION --------------------
\section{Education}
\textbf{Bachelor of Technology in Computer Science and Engineering} \hfill \textbf{Graduated: July 2026} \\
\textit{Indian Institute of Technology Gandhinagar (IITGN)} \hfill \textit{CGPA: 9.59 / 10} \\
\textbf{\href{https://pranavjoshiiitgn.notion.site/54b8560288a346d39574e2b4cd6959ad?v=12b71196f4794d6485295473ab547e0f}{Relevant Coursework}:} Signals \& Systems, Digital Systems, Computer Architecture, Computer Vision, Machine Learning, Data Science, Biostatistics, Computational Physics, Databases.

% -------------------- PUBLICATIONS & POSTERS --------------------
\section{Publications, Posters \& Preprints}
\begin{itemize}
    \item \textbf{Pranav Joshi}, Abhilipsa Sahoo, Joycee Mekie and Himanshu Shekhar, ``Novel Transducer Domain Partitioning Architecture for FPGA-Based Total Focusing Method Ultrasonic Imaging'', \textit{Accepted Poster at IEEE Symposium on Acoustic and Ultrasonic Systems (IEEE SAUS)}, 2026. \textit{(Full manuscript in preparation)}.
    \item \textbf{\href{https://arxiv.org/abs/2507.10019}{Pranav Joshi, ``Sampling-Based Estimation of Jaccard Containment and Similarity''}}, \textit{arXiv preprint arXiv:2507.10019}, 2025.
    \item \textbf{Pranav Joshi} and Ashutosh Srivastava, ``Characterizing Social Jet-Lag Dynamics from Consumer Wearables'', \textit{Manuscript in preparation, Computational Biology Lab (CLOCK), IITGN}, 2026.
\end{itemize}

% -------------------- RESEARCH & PROFESSIONAL EXPERIENCE --------------------
\section{Research \& Professional Experience}

\entryheader{Independent Research: Co-Evolutionary Optimization for Sparse-Array TFM}{Ongoing}{Hardware-Friendly Ultrasonic Beamforming \& Design Space Exploration}{}
\begin{itemize}
    \item Developing a time-complexity optimized Cooperative Coevolutionary Genetic Algorithm (CCGA) to jointly evolve sparse transducer array layouts and non-linear partial-image activation functions.
    \item Replacing resource-heavy $p^{\text{th}}$-root and $p^{\text{th}}$-power coherence operations with learned symbolic and linear regression models over FPGA-implementable basis functions to amplify constructive interference and suppress noise with minimal LUT/DSP overhead.
\end{itemize}

\vspace{4pt}
\entryheader{FPGA Research Contractor}{May -- July 2026}{MUSE Lab, IIT Gandhinagar}{Gandhinagar, India}
\begin{itemize}
    \item Ramped up on ultrasound beamforming theory and background material; documented learnings in a comprehensive \href{https://pranav-joshi-iitgn.github.io/Medical_Imaging_Blog/}{Medical Imaging Systems Blog}
    \item Synthesized and implemented an end-to-end Total Focusing Method (TFM) hardware accelerator for Non-Destructive Testing (NDT) ultrasonic imaging on Nexys 4 DDR and Xilinx AC701 FPGA boards.
    \item Achieved real-time reconstruction at \textbf{30 fps for a $128 \times 128$ pixel grid} on a single FPGA via novel transducer domain partitioning, dynamic programming, deep pipelining, and AXI-stream batching.
    \item Developed a simulation and reconstruction library featuring Software-Assisted Configuration, storing pre-computed delay laws in DDR3 memory to bypass on-chip CORDIC latency and support complex acoustic geometries (e.g., angled wedges).
    \item Engineered a custom link-layer protocol over RMII/RGMII Ethernet to minimize OS context-switching overhead.
    \item Completed technology transfer to industry partner \href{https://modsonic.com/}{Modsonic Instruments Manufacturing Co. Pvt. Ltd.}
\end{itemize}

\vspace{4pt}
\entryheader{Research Intern}{May -- June 2026}{Computational Biology Lab (CLOCK), IIT Gandhinagar}{Gandhinagar, India}
\begin{itemize}
    \item \textbf{Physiological Signal Conditioning:} Proposed a statistical framework to quantify intra-individual Social Jet-Lag (SJL) from high-granularity wearable sensor streams (Ultrahuman). Engineered an unsupervised pipeline using 2D histograms, $k$-means clustering, and unpaired $t$-tests to isolate and filter firmware-induced motion artifacts in Heart Rate (HR) and Heart Rate Variability (HRV).
    \item \textbf{Metric Formulation:} Formulated the Heart-Rate-Autocorrelation (HRA) metric and burstiness index to characterize acute SJL events, controlling for confounding variables via Ordinary Least Squares (OLS) regression and validating superiority over traditional Sleep Midpoint tracking.
    \item \textbf{Time-Series Forecasting:} Implemented a Cosinor autoregressive harmonic regression model to extract circadian phase/amplitude parameters, and benchmarked it against a custom multi-modal Transformer architecture for forecasting physiological trajectories.
\end{itemize}

\vspace{4pt}
\entryheader{Data Engineer (July 2026 -- Present) \& Data Intern (May -- July 2025)}{May 2025 -- Present}{Texas Instruments}{Bengaluru, India}
\begin{itemize}
    \item \textbf{AI-Assisted Hardware System Design (Data Engineer):} Developing AI-driven Block Diagram generation tools for \texttt{ti.com} to automate component selection and accelerate high-level end-equipment architecture design.
    \item \textbf{Large-Scale Statistical Retrieval (Data Intern):} Engineered a metadata-free Retrieval-Augmented Generation (RAG) and multi-agent Text-to-SQL pipeline scaling across thousands of columns and $\sim 10^9$ data points in enterprise Oracle databases. Overcoming rate-limited schema matching directly motivated the mathematical formulation of the arXiv Jaccard Containment estimator.
    \item \textbf{Recognition:} Awarded a Pre-Placement Offer (PPO) for outstanding performance during the internship, leading to the current Data Engineer role.
\end{itemize}

% -------------------- SELECTED ACADEMIC PROJECTS --------------------
\section{Selected Academic Projects}

\projectheader{\href{https://pranav-joshi-iitgn.github.io/PH509/}{Numerical Modeling \& Simulation of Wave and Dynamical Systems}}{PH509, IITGN}{2026}
\begin{itemize}
    \item Developed $\mathcal{O}(N)$ and $\mathcal{O}(MN)$ Time-Dependent Schr\"{o}dinger Equation (TDSE) wave-propagation solvers using Split-Operator and Alternating Direction Implicit (ADI) tridiagonal matrix methods.
    \item Performed von Neumann stability analysis on Space-Discretized Leap-Frog (SDLF) schemes to prevent numerical divergence.
    \item Built adaptive Position-Verlet Leapfrog integrators for 2-body orbital mechanics, verifying conservation laws.
\end{itemize}

\vspace{3pt}
\projectheader{\href{https://github.com/pranav-joshi-iitgn/YajatProject}{Person Velocity Estimation via Projective Geometry \& Computer Vision}}{IITGN}{2024}
\begin{itemize}
    \item Integrated YOLO-based object and head tracking pipelines with native projective geometry homography transformations in OpenCV to estimate real-time physical velocity from monocular video feeds.
\end{itemize}

\vspace{3pt}
\projectheader{\href{https://github.com/pranav-joshi-iitgn/ML_Journey}{Automatic Differentiation Engine \& Classical ML from Scratch}}{ES335, IITGN}{2023}
\begin{itemize}
    \item Engineered an automatic differentiation (autograd) engine supporting dynamic computational graphs, backpropagation, and modular neural network layers. 
    \item Built custom decision tree and random forest ensembles driven by Shannon Entropy and RMSE minimization heuristics.
\end{itemize}

\vspace{3pt}
\projectheader{\href{https://github.com/pranav-joshi-iitgn/Group}{Open Source Contributions to Computational Group Theory}}{IITGN}{Jan -- Jun 2024}
\begin{itemize}
    \item Collaborated with the author of two research papers on group-theoretic algorithms to implement algebraic routines in the GAP and SageMath symbolic computation ecosystems.
\end{itemize}

\vspace{3pt}
\projectheader{Quantized Auto-Encoder Neural Network on FPGA}{Digital Systems, IITGN}{2023}
\begin{itemize}
    \item Designed and synthesized a feed-forward quantized auto-encoder neural network at the Register-Transfer Level (RTL) using SystemVerilog on an FPGA development board.
\end{itemize}

\vspace{3pt}
\projectheader{\href{https://github.com/pranav-joshi-iitgn/AlgoLang}{AlgoLang: Custom Programming Language}}{IITGN}{Jan -- Apr 2025}
\begin{itemize}
    \item Engineered an end-to-end language toolchain comprising a recursive descent parser, lexical analyzer, and AST interpreter in Python for a dynamically typed, lexically scoped programming language.
    \item Built a back-end compiler and custom assembler targeting the \textbf{MIPS32 Instruction Set Architecture (ISA)}, translating high-level control flow and arithmetic expressions into low-level machine code.
    \item Implemented a lightweight Virtual Machine (VM) in C to execute compiled MIPS32 bytecode with custom memory management, register-file emulation, and an interactive shell interface.
\end{itemize}

% -------------------- AWARDS & ACHIEVEMENTS --------------------
\section{Academic Awards \& Achievements}
\begin{itemize}
    \item \textbf{HackRush'25 Hackathon:} Winner in the Computer Architecture track (held at IITGN).
    \item \textbf{Academic Excellence:} Dean's List for semesters 2, 4, 5, and 6.
\end{itemize}

% -------------------- TECHNICAL SKILLS --------------------
\section{Technical Skills}
\begin{itemize}[leftmargin=0pt, label={}]
    \item \textbf{Hardware \& Systems:} Verilog, SystemVerilog, C, Vivado, AXI4, DDR3 (MIG), UART, Ethernet (RMII/RGMII), Timing Analysis, Clock Domain Crossing (CDC), BRAM.
    \item \textbf{Scientific Computing \& ML:} Python, MATLAB, PyTorch, SciPy, NumPy, Pandas, OpenCV, SQL.
    \item \textbf{Software Tools \& Publishing:} LaTeX, Pandoc Markdown, Jupyter, Git, Bash, PowerShell, SSH.
\end{itemize}

\end{document}